\begin{claim}{D.49}{Proved}{The standard centralizer and the two-element conjugator fibres}
Let $P$ be the full pure three-cusp group of D.32 and
$\delta(x)=x^{-1}$, $\delta(y)=y^{-1}$.
Its full centralizer and normalizer are
\[
             C_P(\delta)=N_P(\langle\delta\rangle)
                         =\langle\delta\rangle.         \tag{D.49a}
\]
Define the actual conjugacy orbits
\[
 \mathcal O_P=\{\beta\delta\beta^{-1}:\beta\in P\},
 \qquad
 \mathcal O_{GT_0}=\{\beta\delta\beta^{-1}:\beta\in GT_0\}.
\]
Every fibre of $P\to\mathcal O_P$ is exactly the
right coset $\{\beta,\beta\delta\}$, and
\[
                    \mathcal O_{GT_0}
                            =\mathcal O_P\cap GT_0.     \tag{D.49b}
\]
For $\alpha\in\mathcal O_P$, its two conjugators
have the exact normalized formulas
\[
 \begin{array}{ll}
  x\mapsto x^\mu,& y\mapsto h y^\mu h^{-1},\\
  x\mapsto x^{-\mu},& y\mapsto h y^{-\mu}h^{-1},
 \end{array}
       \qquad \mu\in\widehat\Z^\times,\quad h\in F'.   \tag{D.49c}
\]
Here $h$ is unique, and the unordered pair
$\{\mu,-\mu\}$ is unique.
In D.35's solution space, a choice $(\mu,h,u)$
therefore gives all the solutions over $\alpha$ as
\[
 (\mu,h,u(xy)^a),\qquad
 (-\mu,h,x^{-\mu}u(xy)^a),\qquad a\in\widehat\Z.
                                                        \tag{D.49d}
\]
These assertions concern the specified orbit map;
they do not assert that it contains every full
pure or $GT_0$ involution.
\end{claim}
\begin{proof}
Lochak--Schneps prove that complex conjugation is
self-centralizing in the usual $GT$
\cite[Proposition 4(ii), with Lemma 9]{R-LS}.
Their argument uses the first two defining relations
and applies to $GT_0$ as well.  We give the passage
from the full group $P$ explicitly, together with
the normalized fibre formulas.

Let $\beta\in C_P(\delta)$ have normalized pair
$(\mu,k)$, so
$\beta(x)=x^\mu$ and $\beta(y)=k^{-1}y^\mu k$,
with $k\in F'$.
D.32's normalization is a homomorphic section
into $\Aut(F)$, so the commutation is an actual
automorphism equality.
Its value on $y$ makes $k\delta(k)^{-1}$
centralize $y^\mu$.
The peripheral centralizer is $\langle y\rangle$
by A.7, while $k\delta(k)^{-1}\in F'$.
Abelianization therefore gives $k=\delta(k)$.
The exact free-product identification of D.35,
$F\rtimes_\delta C_2\cong C_2\amalg C_2\amalg C_2$,
gives $C_F(\delta)=1$.  Hence $k=1$.

We also need the precise inclusion in $GT_0$.
For each geometric cusp permutation $\sigma\in\Sigma$,
form $r_\sigma=\beta^{-1}\sigma\beta\sigma^{-1}$.
It lies in $P^1=\ker\lambda$, and commutes with
$\delta$, because both $\beta$ and $\sigma$ do.
D.35 proved $C_{P^1}(\delta)=1$, so $r_\sigma=1$.
The exact identity $GT_0=C_P(\Sigma)$ of D.35
therefore puts this $\beta$ in $GT_0$.

Its second Harbater--Schneps relation, with
$k=1$ and $m=(\mu-1)/2\in\widehat\Z$, is
\[
                (xy)^{-m}y^m x^m=1,
       \qquad\text{equivalently}\quad
                         (xy)^m=y^m x^m.                \tag{D.49e}
\]
For every positive integer $n$ use the explicit
finite Heisenberg quotient with coordinates
$(a,b,c)\in(\Z/n)^3$ and multiplication
\[
 (a,b,c)(a',b',c')=(a+a',b+b',c+c'+ab'),
       \quad x=(1,0,0),\quad y=(0,1,0).
\]
The two sides of (D.49e) have coordinates
\[
 (m,m,m(m+1)/2),\qquad(m,m,0),
                                                        \tag{D.49f}
\]
where the binomial expression is the continuous
integer-valued function on $\widehat\Z$.
These identities first follow for nonnegative
integers by multiplication, and then for all
$m\in\widehat\Z$ by continuity; modulo $n$ they
depend only on $m\bmod2n$.
Equality in every such quotient gives
$m(m+1)/2=0$ in $\widehat\Z$, and hence
\[
                           \mu^2=1.                     \tag{D.49g}
\]
Since $k=1$, the actual square of $\beta$ fixes
both $x$ and $y$ by (D.49g).  Thus $\beta^2=1$.
D.34 gives the global multiplier $-1$ for a
nonidentity pure involution; the identity has
multiplier $1$.
The normalized pair is therefore $(1,1)$ or
$(-1,1)$, proving $C_P(\delta)=\langle\delta\rangle$.
A normalizer of its order-two subgroup must
centralize its sole nonidentity element, proving
the normalizer equality in (D.49a).

Two elements $\beta,\beta'$ have the same image
under the orbit map exactly when
$\beta^{-1}\beta'\in C_P(\delta)$.
This proves its two-element fibres.
For an orbit element lying in $GT_0$, D.35's
proved equality of conjugator sets places both
conjugators in $GT_0$; the reverse inclusion is
immediate.  This proves (D.49b).

Write one conjugator in D.35's convention as
$\beta_{\mu,h}$.  Right composition with $\delta$
gives exactly $\beta_{-\mu,h}$, so (D.49c) follows
with the same $h$.
They are distinct because $\mu\ne-\mu$ in
$\widehat\Z^\times$.
Moreover actual normalized conjugation identifies
$C_F(\alpha)$ with $\beta(C_F(\delta))=1$.
The description of all coboundaries in (D.35g)
then makes $h$ the unique solution of
$h^{-1}\alpha(h)=f$ in $F'$.
The two-element fibre proves uniqueness of the
unordered multiplier pair.

Finally put $A=x^\mu$, $B=h y^\mu h^{-1}$ and
suppose $AB=u(xy)^\mu u^{-1}$.
The exact word identity
\[
           A^{-1}B^{-1}=A^{-1}(AB)^{-1}A
\]
shows that $x^{-\mu}u$ is a third-cusp conjugator
for the opposite multiplier.
D.35's centralizer calculation gives all other
third-cusp conjugators by right multiplication
by $(xy)^a$.  This proves (D.49d).
The program
\path{code/part_a/standard_involution_centralizer.py}
checks (D.49f) over sixteen finite moduli and
checks the displayed word identity.
The full conclusion (D.49a) uses the proved
global-sign theorem D.34, beyond those finite
checks.
\end{proof}
